\subsection{三角函数线}

\begin{definition}[任意角的三角函数定义]
    设角 $\alpha$ 的顶点为坐标原点，始边为 $x$ 轴的非负半轴，终边上任意一点 $P$ 的坐标为 $(x, y)$，它与原点的距离为 $r = \sqrt{x^2+y^2}$ ($r>0$)，则：
    $$\sin\alpha = \frac{y}{r}, \quad \cos\alpha = \frac{x}{r}, \quad \tan\alpha = \frac{y}{x}$$
    \textbf{全是天才：一句话记忆象限}
\end{definition}



\begin{figure}[H]
    \centering
    \begin{tikzpicture}[scale=1]
        \draw (0,0) circle (2);
        \coordinate (origo) at (0,0);
        \coordinate[label = above left:$P$] (P) at (1,{sqrt(3)});
        \draw[thick,->] (origo) -- ++(3,0) node[black,right] (xaxis){$x$};
        \draw[thick,gray,->] (origo) -- ++(0,3) node (mary) [black,above] {$y$};
        \draw[thick] (origo) -- ++(60:2) coordinate (bob);
        \pic [draw, ->, red, angle eccentricity=2] {angle = xaxis--origo--bob} node[left, color=red] {$\theta$};
        \draw[blue, ->] (1,{sqrt(3)}) -- (1,0);
        \draw[cyan, ->] (1,{sqrt(3)}) -- (0,{sqrt(3)});
        \node[color=cyan] at (-0.5,{sqrt(3)}) {$\cos \theta$};
        \node[color=blue] at (1,-0.3) {$\sin \theta$};
        \draw[black,dotted] (-2,2)--(3,2);
        \draw[black,dotted] ({2*sqrt(3)/3},2)--(1,{sqrt(3)});
        \draw[brown,->] ({2*sqrt(3)/3},2)--({2*sqrt(3)/3},0);
        \node[color=brown] at (1.6,1) {$\tan \theta$};
        \node at (P)[circle,fill,inner sep=2pt]{};
    \end{tikzpicture}
    \caption{三角函数线}
\end{figure}

\begin{example}
    已知角 $\alpha$ 的终边经过点 $P(-3, 4)$，求 $\sin\alpha$, $\cos\alpha$, $\tan\alpha$ 的值。
\end{example}
\begin{solution}
    点 $P(-3, 4)$ 中，$x=-3, y=4$，所以 $r = \sqrt{(-3)^2 + 4^2} = \sqrt{9+16} = \sqrt{25} = 5$。
    $$\sin\alpha = \frac{y}{r} = \frac{4}{5}, \quad \cos\alpha = \frac{x}{r} = -\frac{3}{5}, \quad \tan\alpha = \frac{y}{x} = -\frac{4}{3}$$
\end{solution}
\vspace{2em}

\subsection{齐次化切问题}

\begin{theorem}[同角三角函数基本关系]
    \begin{itemize}
        \item \textbf{平方关系:} $\sin^2\alpha + \cos^2\alpha = 1$
        \item \textbf{商数关系:} $\tan\alpha = \frac{\sin\alpha}{\cos\alpha}$
    \end{itemize}
\end{theorem}

\begin{example}
    已知 $\alpha$ 是第三象限角，且 $\cos\alpha = -\frac{12}{13}$，求 $\sin\alpha$ 和 $\tan\alpha$ 的值。
\end{example}
\begin{solution}
    因为 $\alpha$ 是第三象限角，所以 $\sin\alpha < 0$ 且 $\tan\alpha > 0$。\\
    根据平方关系 $\sin^2\alpha + \cos^2\alpha = 1$，有：
    $\sin^2\alpha = 1 - (-\frac{12}{13})^2 = 1 - \frac{144}{169} = \frac{25}{169}$。\\
    因为 $\sin\alpha < 0$，所以 $\sin\alpha = -\frac{5}{13}$。\\
    根据商数关系：
    $\tan\alpha = \frac{\sin\alpha}{\cos\alpha} = \frac{-5/13}{-12/13} = \frac{5}{12}$。
\end{solution}

\begin{theorem}
    齐次化切通常在分式的上下同时除以 $cosx$ 的某次方，以将 $sinx$ 和 $cosx$ 的幂次化为相同的次数，从而将分式化为一个关于 $\tan x$ 的有理函数。
\end{theorem}

\begin{example}
    $$\text{已知：}\frac{\sin (2 \pi+\theta) \tan (\pi+\theta) \tan (3 \pi-\theta)}{\sin (\pi-\theta) \tan (-\pi-\theta)}=1 \text{，则} \sin ^2 \theta+3 \sin \theta \cos \theta+2 \cos ^2 \theta \text{的值是（\h{2em}）}$$
\end{example}

\v{8em}

\begin{example}
    已知 $\tan (\pi+\alpha)=-2$
    \begin{enumerate}
        \item 求 $\sin\alpha$ 和 $\cos\alpha$ 的值；
        \item 求 $\sin^2\alpha + 3\sin\alpha\cos\alpha + 2\cos^2\alpha$ 的值。
    \end{enumerate}
\end{example}

\v{8em}

\begin{theorem}[诱导公式]
    奇变偶不变，符号看象限
\end{theorem}

\begin{example}
    计算 $\sin(-600^\circ) + \cos(\frac{11\pi}{6})$ 的值。
\end{example}
\begin{solution}
    $\sqrt{3}$。
\end{solution}
\vspace{2em}

\begin{problem}
$$\text{已知：}f(\beta)=\frac{\sin (\pi-\beta) \cos (2 \pi-\beta) \tan (\beta+\pi)}{\tan (-\beta-\pi) \sin (-\pi-\beta)}$$
\begin{itemize}
    \item （1）若角 $\beta$ 是第三象限角，且 $\sin (\beta-\pi)=\frac{1}{5}$ ，求 $f(\beta)$ 的值；
    \item （2）若 $\beta=\frac{37 \pi}{3}$ ，求 $f(\beta)$ 的值．
\end{itemize}
\end{problem}

\v{8em}

\begin{problem}
化简:
$$
    \frac{\tan (\pi+\alpha) \cos (2 \pi+\alpha) \sin \left(\alpha-\frac{3 \pi}{2}\right)}{\cos (-\alpha-3 \pi) \sin (-3 \pi-\alpha)}
$$
\end{problem}


\subsection{凑角问题}

\begin{theorem}[两角和与差的三角函数]
    $$ \sin(\alpha \pm \beta) = \sin\alpha\cos\beta \pm \cos\alpha\sin\beta $$
    $$ \cos(\alpha \pm \beta) = \cos\alpha\cos\beta \mp \sin\alpha\sin\beta $$
    $$ \tan(\alpha \pm \beta) = \frac{\tan\alpha \pm \tan\beta}{1 \mp \tan\alpha\tan\beta} $$
\end{theorem}

\begin{example}
    已知 $\cos\alpha = \frac{1}{7}$, $\cos(\alpha+\beta)=-\frac{11}{14}$，且 $\alpha, \beta$ 均为锐角，求 $\cos\beta$ 的值。
\end{example}
\begin{solution}
    因为 $\alpha, \beta$ 均为锐角，所以 $\sin\alpha > 0$ 且 $\sin\beta > 0$。\\
    由 $\cos\alpha = \frac{1}{7}$，得 $\sin\alpha = \sqrt{1 - (\frac{1}{7})^2} = \frac{4\sqrt{3}}{7}$。\\
    又因为 $\alpha+\beta \in (0, \pi)$，且 $\cos(\alpha+\beta)=-\frac{11}{14} < 0$，所以 $\alpha+\beta$ 是钝角。\\
    $\sin(\alpha+\beta) = \sqrt{1 - (-\frac{11}{14})^2} = \sqrt{\frac{75}{196}} = \frac{5\sqrt{3}}{14}$。\\
    $\cos\beta = \cos((\alpha+\beta)-\alpha) = \cos(\alpha+\beta)\cos\alpha + \sin(\alpha+\beta)\sin\alpha$
    $= (-\frac{11}{14})(\frac{1}{7}) + (\frac{5\sqrt{3}}{14})(\frac{4\sqrt{3}}{7}) = -\frac{11}{98} + \frac{60}{98} = \frac{49}{98} = \frac{1}{2}$。
\end{solution}
\vspace{2em}


\begin{theorem}[二倍角公式与降幂公式]
    $$ \sin(2\alpha) = 2\sin\alpha\cos\alpha \quad \text{（交叉项降幂）}$$
    $$ \cos(2\alpha) = \cos^2\alpha - \sin^2\alpha = 2\cos^2\alpha - 1 = 1 - 2\sin^2\alpha $$
    $$ \sin^2\alpha = \frac{1-\cos(2\alpha)}{2}, \quad \cos^2\alpha = \frac{1+\cos(2\alpha)}{2} \quad \text{（平方降幂）}$$
\end{theorem}

\begin{example}
    已知函数 $f(x)=\sin x+\sin ^2 \frac{x}{2}-\cos ^2 \frac{x}{2}$.求函数 $f(x)$ 的单调递增区间；
\end{example}

\vspace{2em}

\begin{theorem}[辅助角公式]
    用于将 $a\sin x + b\cos x$ 的形式化为单一函数：
    $$ a\sin x + b\cos x = \pm \sqrt{a^2+b^2} \sin(x + \varphi) $$
    其中 $\tan\varphi = \frac{\text{cos系数}}{\text{sin系数}}$。

    注意要点：
    \begin{enumerate}
        \item 整个式子的符号和 $sin$ 的正负一致。
        \item $\varphi$ 的取值范围是 $(-\frac{\pi}{2}, \frac{\pi}{2})$。
    \end{enumerate}
\end{theorem}

\begin{example}
    求函数 $f(x) = \sin x - \sqrt{3}\cos x$ 在区间 $[0, \pi]$ 上的最大值和最小值。
\end{example}

\v{8em}

\begin{example}
    已知 $\overrightarrow{m}=\left(\sin 2 x, \sin ^2 x\right), \overrightarrow{n}=(1,2), f(x)=\overrightarrow{m} \cdot \overrightarrow{n}$ ．
    \begin{enumerate}
        \item 求 $f(x)$ 的最小正周期；
    \end{enumerate}
\end{example}

\v{8em}

\begin{example}
    求下列函数的值域．
    \begin{enumerate}
        \item $f(x)=\cos x(\sin x+\sqrt{3} \cos x)-\frac{\sqrt{3}}{2} \quad\left(x \in\left[0, \frac{\pi}{2}\right]\right)$ ；
        \item $y=\cos ^2 x+\sin x-1$ ；
        \item $f(x)=\sin x \cos x+\sqrt{2} \sin \left(x-\frac{\pi}{4}\right)$.
    \end{enumerate}
\end{example}
\v{8em}

\vspace{2em}